\SourceLeaf{231}%
Ich habe hier bloß die Grundzüge des Beweises angeben wollen, ohne jede Feinarbeit oder Politur, für die ich jetzt keine Zeit habe. Es steht Dir frei, es Deinem Sohn mitzuteilen. Jedenfalls bitte ich Dich, ihn herzlich von mir zu grüßen und ihm meine besondere Hochachtung zu versichern. Fordere ihn aber doch zugleich auf, sich mit der folgenden Aufgabe zu beschäftigen.

»Den Kubikinhalt des Tetraeders (von vier Ebenen begrenzten Raumes) zu bestimmen«.

Da der Flächeninhalt eines Dreiecks sich so einfach angeben lässt, hätte man erwarten sollen, dass es auch für diesen Kubikinhalt einen ebenso einfachen Ausdruck geben werde. Aber diese Erwartung wird, wie es scheint, enttäuscht.

Um die Geometrie von Anfang an ordentlich zu behandeln, ist es unerlässlich, die Möglichkeit eines Planums zu beweisen. Die gewöhnliche Definition enthält zu viel und impliziert eigentlich subreptiv schon ein Theorem. Man muss sich wundern, dass alle Schriftsteller von Euklid bis in die neueste Zeit dabei so nachlässig vorgegangen sind. Diese Schwierigkeit ist jedoch von ganz anderer [Natur] als die Schwierigkeit, zwischen $\Sigma$ und $S$ zu entscheiden, und die erstere ist nicht allzu schwer zu beheben. Wahrscheinlich wird mich auch schon Dein Buch in dieser Frage zufriedenstellen.

Gerade in der Unmöglichkeit, zwischen $\Sigma$ und $S$ a priori zu entscheiden, liegt der klarste Beweis, dass Kant Unrecht hatte zu behaupten, der Raum sei nur \textls{Form} unserer Anschauung. Einen anderen ebenso starken Grund habe ich in einem kleinen Aufsatz angedeutet, der in den Göttingischen Gelehrten Anzeigen 1831 steht, Stück 64, pag.~625. Vielleicht wird es Dich nicht gereuen, wenn Du Dich bemühst, Dir diesen Band der G.G.A. zu verschaffen (was jeder Buchhändler in Wien oder Ofen leicht bewirken kann), da darin unter anderem auch die Quintessenz meiner Ansicht von den imaginären Größen auf ein paar Seiten dargelegt ist. \mbox{.\ .\ .\ .\ .\ .\ .}

